\chapter{Background}\label{chap:background}

This chapter covers the background that the rest of the thesis builds on. Section~\ref{sec:bg-tcp-cc} goes over the early TCP implementation and the first congestion control measures Section~\ref{sec:bg-quic} introduces the QUIC protocol and some of the features that distinguish it from TCP and Section~\ref{sec:bg-traffic} summarizes measurements of Internet traffic.

\section{TCP and the initial congestion control solutions}\label{sec:bg-tcp-cc}

\subsection{Congestion collapse and the first congestion control}\label{sec:bg-origins}

In a packet-switched network, links are shared, and networks do not stop
senders from uploading more traffic than a link can carry. The excess packets that are generated wait in
router buffers, and once the buffers are full, they are dropped, leading to packet loss. In 1984, Nagle described how this can turn into congestion collapse~\cite{rfc896}. A sudden load
raises the round-trip time (RTT) faster than hosts can update their estimates of
it, so hosts retransmit packets that are only delayed. The buffers fill with
duplicate copies, every packet is sent several times, and throughput falls to a
small fraction of normal.

The problem became real in 1986, with the first of a series of congestion
collapses on the early Internet: throughput between
the Lawrence Berkeley Laboratory and UC~Berkeley dropped from 32\,kbit/s to 40\,bit/s. 
In an attempt to mitigate the decrease in throughput, seven algorithms were added to TCP, among them round-trip-time variance estimation, exponential retransmit timer backoff, slow start and dynamic window sizing on congestion~\cite{jacobson1988congestion}.

These algorithms focus on maintaining the conservation of packets: a
connection running with a full window should put a new packet
into the network only when an old one leaves it. Acknowledgments (ACKs) act
as a clock, because each ACK reports that a packet has left the network and
allows the sender to send another. Jacobson identified three ways in which
conservation can fail: the connection never reaches the full window, the sender
injects a new packet before an old one has left, or resource limits along the path
prevent equilibrium altogether~\cite{jacobson1988congestion}. Each of the following algorithms addresses one of them.

\paragraph{Slow start.} Slow start adds a congestion window (\textit{cwnd}) to the
connection state. The sender sets \textit{cwnd} to one packet when it starts or
restarts after a loss, increases it by one packet for every ACK of new data, and
never has more than the minimum of \textit{cwnd} and the receiver's advertised
window (RCV.WND) in flight. Each ACK therefore allows for the release of two packets, so the \textit{cwnd}
doubles every round-trip-time (RTT): reaching a window of $W$ packets takes $RTT\log_2 W$.

\paragraph{Retransmission timer.} The second failure is a retransmission timer that
fires while a packet is still in transit. The original TCP specification~\cite{rfc793} set the
timeout to a fixed multiple of a smoothed RTT estimate. Jacobson
showed that the current specification tolerated loads of up to about
30\% of the capacity: above that, the RTT varies so much that the timer retransmits packets that
are just delayed, adding more load to a network that is already
overloaded~\cite{jacobson1988congestion}. His new timer estimates the variation of the
RTT as well as its mean, sets the timeout from both, and backs off exponentially
when retransmissions repeat. Together with Karn's rule~\cite{karn1987improving}, which discards RTT samples
from retransmitted packets, this design is still the basis
of TCP's retransmission timer, now standardized in RFC 6298~\cite{rfc6298}.

\paragraph{Congestion avoidance.} To mitigate the third failure, the network should signal congestion when the path does not have enough capacity and the endpoints must react to it. Jacobson observed that on most paths packet loss due to damage is
rare, well below 1\%, so a lost packet almost always means that a buffer
overflowed~\cite{jacobson1988congestion}. Packet loss was therefore chosen as the congestion signal, requiring no changes to routers. Congestion avoidance reacts to packet loss with
additive increase, multiplicative decrease (AIMD): on a timeout,
\textit{cwnd} is halved, and on each ACK of new data, \textit{cwnd} grows by
$1/\textit{cwnd}$, about one packet per RTT. On a timeout, half of the current window is stored in a slow start
threshold (\textit{ssthresh}) and \textit{cwnd} restarts at one packet. Slow start
then runs until \textit{cwnd} reaches \textit{ssthresh}, and congestion avoidance
takes over above it.

Together with \emph{fast retransmit}, which retransmits a lost packet after a small
number of duplicate ACKs instead of waiting for the timer, these algorithms were
first released in 4.3BSD Tahoe in 1988 and are therefore known as TCP
Tahoe~\cite{fall1996simulation}. Their effect is still relevant today: the timer improvements
and congestion control mechanisms of TCP have largely removed
congestion collapse~\cite{rfc2914}, and end-host control with packet
loss as the congestion signal is still the basis of the default TCP algorithms.

\subsection{CUBIC}\label{sec:bg-cubic}

\paragraph{Fast recovery} TCP Reno~\cite{rfc5681} kept Tahoe's algorithms and added \emph{fast recovery}.
When three duplicate ACKs indicate that a single packet was lost, the sender
retransmits it, halves \textit{cwnd} and continues in congestion avoidance, because
the duplicate ACKs show that later packets are still leaving the network. Only a
retransmission timeout sends the connection back to slow start with a window of one
packet. In steady state, the window therefore follows a sawtooth: it grows by one
packet per RTT and is halved at each loss. This pattern can also be seen in Figure~\ref{fig:tcp_reno}

\begin{figure}[ht]
    \centering
    \includegraphics[width=1\linewidth]{figures/TCP_Reno_Sawtooth.png}
    \caption{TCP Reno in steady state (from Computer Networks: A Systems Approach~\cite{peterson2022computer}) }
    \label{fig:tcp_reno}
\end{figure}

This sawtooth does not scale well to paths with a large bandwidth-delay product (BDP).
After a loss, Reno needs cwnd/2 RTTs to recover. On a 1\,Gbit/s
path with an RTT of 100\,ms, the window that fills the path is about 8\,300 packets
of 1\,500 bytes. After a single loss, Reno needs about 4\,150 RTTs, or seven
minutes, to return to it. This problem is noticeable even with
windows of a few hundred packets. Several variants of TCP
were proposed to solve it; the one that was eventually widely adopted is CUBIC.

\paragraph{The cubic window function.} Ha, Rhee and Xu proposed CUBIC~\cite{ha2008cubic} in 2008 as the
successor to BIC-TCP. It was published as RFC 9438
in 2023~\cite{rfc9438}. CUBIC keeps Reno's structure of slow start, congestion
avoidance and multiplicative decrease, but changes two of its rules. First, on a
congestion event, it remembers the current window as $W_{\max}$ and reduces the
window by a factor of $\beta = 0.7$ instead of $0.5$. Second, in congestion
avoidance, it computes the window from a cubic function of the time $t$ since the
start of congestion avoidance:
\begin{equation}
  W_{\mathit{cubic}}(t) = C\,(t-K)^3 + W_{\max},
  \qquad
  K = \sqrt[3]{\frac{W_{\max} - \mathit{cwnd}_{\mathit{epoch}}}{C}},
  \label{eq:bg-cubic}
\end{equation}
where $\mathit{cwnd}_{\mathit{epoch}}$ is the window after the reduction, windows
are measured in segments, $t$ and $K$ are in seconds, and $C = 0.4$. $K$ is the time
the function takes to return to $W_{\max}$. The window first grows quickly, then
slows down as it approaches $W_{\max}$ (the concave region) and stays almost flat
around it, because that is where the last loss happened. If no loss occurs there,
the window grows past $W_{\max}$, slowly at first and then faster (the convex
region), to probe for capacity that has become available. The cwnd growth can also be observed in Figure~\ref{fig:tcp_cubic}.

\begin{figure}[ht]
    \centering
    \includegraphics[width=0.6\linewidth]{figures/TCP_Cubic.png}
    \caption{Cwnd as a function of time after loss for CUBIC (from Computer Networks: A Systems Approach~\cite{peterson2022computer})}
    \label{fig:tcp_cubic}
\end{figure}

CUBIC changes congestion
avoidance and the response to loss, but not the starting phase of a connection: while
\textit{cwnd} is at most \textit{ssthresh}, it uses a slow start
algorithm, so its window grows from the initial window as in Reno:
RFC~9438 recommends HyStart++~\cite{rfc9406} for this phase.

\paragraph{Deployment.} CUBIC replaced BIC-TCP as the default congestion control of
Linux in kernel 2.6.19~\cite{hemminger2006cubic}, and it is still
the default today. RFC~9438 reports that it is also the
default in the Windows and Apple stacks and calls it the currently most widely
deployed standard for TCP congestion control~\cite{rfc9438}. The RFC also allows
CUBIC to be used by other transport protocols, including QUIC.

\subsection{The initial window}\label{sec:bg-iw}

While the slow start algorithm determines how fast the congestion window can grow, the initial window
(IW) determines at what value the congestion window starts. For short flows, the IW set determines the number of RTTs needed for a transfer to complete, and it has been raised several times since 1988.

Jacobson's slow start began with a window of one packet~\cite{jacobson1988congestion},
and RFC~2001 specified an IW of one segment~\cite{rfc2001}. RFC~2581 raised the
standard IW to two segments and allowed a larger value as an
option~\cite{rfc2581}. That larger value, first proposed as an experiment in RFC~2414
and standardized in RFC~3390 in 2002, is
three segments, or about 4\,KB~\cite{rfc2414,rfc3390}. Deployment followed slowly: in measurements of web servers published in 2005, most servers used an IW of one or two segments, and fewer used three or four~\cite{medina2005measuring}.

In 2010, Dukkipati et al.~\cite{dukkipati2010argument} argued that the IW had not changed since 2002, while
access speeds had increased greatly, and proposed raising it to at least ten
segments, about 15\,KB. In large-scale experiments on
Google's servers, an IW of ten segments improved the average latency of HTTP
responses by about 10\%, with the largest gains on paths with a high RTT or BDP,
and increased the average retransmission rate by 0.5\%.
Linux adopted this IW in kernel 2.6.39 and RFC~6928~\cite{rfc6928} specified it in 2013.

In practice, the IW is not a single value. A scan of the IPv4 address space in
2017 found that IWs of 1, 2, 4 and 10 segments together account for more than 97\% of
web servers, and that among the most popular websites, more than 80\% use an IW of
ten~\cite{ruth2017large}. Content delivery networks go further: they tune the IW per
network, customer or content, mostly below 50 segments, but up to 100
segments~\cite{ruth2019iw}. 

\subsection{Active queue management and ECN}\label{sec:bg-ecn-aqm}

Loss-based congestion control learns about congestion only when a router's buffer
overflows. A router with a simple first-in, first-out queue drops packets only once
its buffer is full, so until then, every packet waits behind the whole queue. Large
buffers at the edge of the network therefore add long delays, a problem known as
bufferbloat~\cite{gettys2012bufferbloat}.

\paragraph{Active queue management.} With active queue management (AQM), a router
drops or marks packets before its buffer is full, to keep the queue short while still
signaling congestion to the senders. The IETF recommends its widespread
deployment~\cite{rfc7567}. CoDel~\cite{rfc8289}, for example, controls the time that packets spend in
the queue rather than the length of the queue: it drops or marks packets when the
minimum queueing delay over an interval, 100\,ms by default, stays above a target of
5\,ms. FQ-CoDel~\cite{rfc8290} also hashes flows into separate queues, serves them in
round robin and applies CoDel to each, and it serves flows that do not build a
queue, such as new and sparse flows, first.

\paragraph{Explicit Congestion Notification.} ECN~\cite{rfc3168} lets a router signal congestion
without dropping packets. Two bits in the IP header carry four
codepoints: a sender marks its packets as ECN-capable with ECT(0) or ECT(1), and a
congested router with AQM changes the codepoint to Congestion Experienced (CE)
instead of dropping the packet. The receiver echoes the mark to the sender, which
reacts to the congestion signal without losing the packet and
without waiting to detect a loss. The original specification requires the response
to CE to be the same as the response to a loss, however RFC~8311~\cite{rfc8311} relaxed this
requirement to allow experiments, for example with a smaller reduction for marks and
with new uses of ECT(1).

\paragraph{Deployment.} Both mechanisms depend on the network. ECN only helps where
the bottleneck runs an AQM that marks packets and where no device on the path clears
the ECN field. In measurements from 2021 and 2022, one cloud provider and two cellular
providers cleared, or ``bleached'', the ECN field as a matter of policy~\cite{lim2022fresh}. 

% \subsection{Metrics for assessing congestion control algorithms}\label{sec:bg-metrics}

% A congestion control algorithm serves two parties. The network operator wants the
% bottleneck used efficiently and shared fairly among flows, while a user wants
% their own transfer to finish quickly.

% \paragraph{Flow completion time.} The flow completion time (FCT) of a transfer is the
% time from its start until its last byte has been received. Dukkipati and McKeown~\cite{dukkipati2006why}
% argue that it is arguably the most important metric for users, who want their
% transfers to complete as quickly as possible, whereas congestion control research
% has focused almost entirely on throughput, utilization and fairness, which matter
% more to the operator. With typical Internet flow sizes,
% they show that TCP Reno makes flows last much longer than necessary, often by one
% or two orders of magnitude. 

% \paragraph{Fairness.} Fairness describes
% how the capacity of a bottleneck is divided among the flows that share it. Jain, Chiu and Hawe defined a fairness index~\cite{jain1984quantitative} for $n$ flows
% with throughputs $x_1, \dots, x_n$:
% \begin{equation}
%   J = \frac{\left(\sum_{i=1}^{n} x_i\right)^2}{n \sum_{i=1}^{n} x_i^2}.
%   \label{eq:jain-index}
% \end{equation}
% The index is independent of the amount of the shared resource. It equals 1 when all
% flows receive the same throughput and falls to $1/n$ when one flow takes all of it.
% Chiu and Jain~\cite{chiu1989analysis} showed that AIMD satisfies the sufficient conditions for converging
% to an efficient and fair allocation regardless of the starting
% state.

% As new algorithms appeared, fairness came to be judged relative to TCP. Floyd and
% Fall~\cite{floyd1999promoting} call a flow \emph{not TCP-friendly} if its long-term arrival rate exceeds that
% of any conformant TCP in the same circumstances. Ware et
% al.\ argue that the bar for deploying a new algorithm should instead be the
% \emph{harm} it causes to existing traffic~\cite{ware2019beyond}. Harm is measured
% against a baseline of the existing traffic running on its own, so a new algorithm
% is not blamed when existing algorithms perform poorly because of their own
% limitations. It also covers any performance metric, including throughput,
% latency, loss rate and flow completion time. Current IETF guidance takes a similar
% view~\cite{rfc9743}. A proposed algorithm must be evaluated against standard IETF
% congestion controls, and the measure of harm is not restricted to the share of
% capacity but also includes latency and loss. The guidance also requires an evaluation of how short flows affect long-lived flows, and of how the new algorithm affects both short-lived and
% long-lived flows that use other common algorithms.


\section{QUIC}\label{sec:bg-quic}

QUIC~\cite{langley2017quic} is a transport protocol that runs over UDP, integrates the TLS~1.3 handshake,
multiplexes independent streams in one connection, and encrypts almost all of its
packets. Google launched an early version as an experiment in
2013 and built it in user space, so that it could be deployed and changed at the pace
of application updates. By 2017, QUIC carried over 30\% of Google's egress traffic,
an estimated 7\% of all Internet traffic. The IETF standardized
QUIC in 2021: RFC~9000 defines the transport~\cite{rfc9000}, RFC~9001 its use of
TLS~\cite{rfc9001}, and RFC~9002 its loss detection and congestion
control~\cite{rfc9002}.

RFC~9002 specifies a congestion controller similar to TCP NewReno, but a sender may choose another algorithm, provided that it follows
the IETF guidelines for congestion control over UDP. Because QUIC runs
in user space, such an algorithm can be deployed together with an application,
without changes to the operating system kernel.

\subsection{Differences between TCP and QUIC}\label{sec:bg-quic-features}

This section describes features of QUIC that are not available in TCP and could be used in congestion control algorithms made specifically for QUIC.

\paragraph{Fewer round trips before the first request.} A TCP connection starts with
a three-way handshake, which adds one RTT before data can be sent, and
TLS~1.3~\cite{rfc8446} adds another round trip for its own handshake. An HTTPS
request over TCP therefore leaves the client after two RTTs. QUIC combines the
transport and cryptographic handshakes, so a new client sends its
request after one RTT. A client that resumes an earlier connection can send its
request in its very first flight, as 0-RTT data. The server may also
include data in its first flight, before the handshake completes, which RFC~9000
calls ``0.5-RTT data''; this does not change the one RTT that a new
client waits before its request leaves. TCP Fast Open~\cite{rfc7413} can also save one RTT, but it
is only experimental. For a short flow, these saved round trips are a
latency gain that no TCP congestion controller can achieve. The handshake also gives the
server an RTT sample before its first data packet, because endpoints must acknowledge
Initial and Handshake packets immediately, so this sample contains no
acknowledgment delay.

\paragraph{Extension frames.} QUIC can be extended with new frame types, which lets
the endpoints exchange information that the base protocol does not carry. QUIC frames
are not self-describing, and an endpoint must treat a frame of unknown type as a
connection error, so an extension first negotiates its frames: each endpoint
announces in a transport parameter which extension frames it is willing to receive,
and endpoints ignore transport parameters that they do not support~\cite{rfc9000}.
Extension frames must be congestion controlled and must elicit an acknowledgment,
unless they replace or supplement the ACK frame. New frame types are registered with
IANA, and provisional registrations are available for experiments.
Several extensions already use this mechanism: the unreliable DATAGRAM
frame~\cite{rfc9221}, the ACK\_FREQUENCY and IMMEDIATE\_ACK frames~\cite{iyengar2026ackfrequency} are all registered with IANA already.

TCP is much harder
to extend. New functionality has to fit into TCP options, which share at most
40~bytes of header space~\cite{rfc793}, and that space is already
scarce. Options are also visible to the network: in
measurements from 2011, only 80\% of the paths tested on port 80 passed TCP options
without interference, against 96\% on an unused port~\cite{honda2011extend}. QUIC
frames travel inside the encrypted and authenticated payload, so middleboxes can
neither read nor modify them.

\paragraph{Path state carried across connections.} A TCP host can cache path state,
such as the RTT and congestion window of earlier connections, but only
locally~\cite{rfc9040}. In QUIC, the server can store state in the TLS session ticket
that the client presents when it resumes a connection~\cite{rfc9001}. RFC~9002 allows a
resumed connection over the same network to start from the previous connection's
smoothed RTT~\cite{rfc9002}. Careful Resume~\cite{rfc9959}, the IETF method for reusing congestion control state, limits
the first RTT of a resumed connection to the initial window and the following jump to
half of the saved congestion window, after which the sender validates the
path.

\paragraph{Packet numbers are never reused.} TCP retransmits lost data with the same
sequence numbers, so an acknowledgment of retransmitted data does not reveal which
transmission arrived. Karn's rule therefore discards RTT samples from retransmitted
segments~\cite{karn1987improving}. A QUIC endpoint must not reuse a packet number
within a packet number space: a lost frame is resent in a new packet
with a new number, so each acknowledgment refers to exactly one
transmission~\cite{rfc9002}. RTT samples are therefore unambiguous, also during loss
recovery, and spurious retransmissions are trivially detected.
A sample is taken only from the largest
newly acknowledged packet of an ACK frame, and it still includes the receiver's
acknowledgment delay and any queueing delay.

\paragraph{Explicit acknowledgment delay.} A TCP receiver may delay an ACK by up to
500\,ms after the first unacknowledged segment arrives~\cite{rfc5681}, and the sender cannot tell by
how much, so the delay inflates its RTT samples. QUIC makes the delay explicit. Each
endpoint announces in its transport parameters the maximum time by which it will
delay its acknowledgments, \texttt{max\_ack\_delay}, which defaults to 25\,ms~\cite{rfc9000}, and
every ACK frame carries an ACK Delay field, the time for which the receiver held the
acknowledgment of the largest acknowledged packet. RFC~9002 subtracts this delay from
the RTT sample, capped at \texttt{max\_ack\_delay} once the handshake is confirmed,
but never lowers a sample below the minimum RTT. The field is self-reported and
covers only the largest acknowledged packet, but it separates the receiver's delay
from the network's. The sender also adds \texttt{max\_ack\_delay} to its probe timeout
(PTO), which in QUIC replaces TCP's retransmission timeout (RTO) and tail loss probe~\cite{rfc9002}.
The PTO follows the RTO computation of RFC~6298, but uses \texttt{max\_ack\_delay}
instead of a fixed minimum:
\begin{equation}
  \mathit{PTO} = \mathit{smoothed\_rtt} + \max(4 \cdot \mathit{rttvar},\
  k_{\mathit{Granularity}}) + \mathit{max\_ack\_delay}.
  \label{eq:bg-pto}
\end{equation}
When the PTO expires without an acknowledgment, the sender sends one or two probe
packets. Unlike an RTO, an expired PTO does not reset the congestion window, because
the loss of a single packet at the end of a flow does not indicate persistent
congestion; QUIC reduces the window to its minimum only when it declares persistent
congestion.

The value of \texttt{max\_ack\_delay} is set by the endpoint that sends the
acknowledgments. In a download, the
client chooses it, and the server cannot change it in the base protocol. A client
that advertises a small value shortens the server's PTO, so a packet lost at the end
of a short flow is recovered sooner, at the cost of more acknowledgments.



\paragraph{Many, irrevocable acknowledgment ranges.} A TCP SACK~\cite{rfc2018} option holds at most
four blocks, or three if timestamps are used, and a receiver may renege on data it has
already acknowledged selectively. A QUIC ACK frame can carry many ACK
ranges, limited only by the packet size, and acknowledgments can never be
revoked~\cite{rfc9002}. The sender therefore knows exactly which packets have
arrived and how many bytes are still in flight. Loss itself is still inferred: a
packet is declared lost when a packet sent at least three packets later is
acknowledged, or after a time threshold of 9/8 of the RTT, and a
declared loss may later turn out to be spurious.

\paragraph{Richer ECN feedback.} With classic ECN, a TCP receiver echoes congestion marks with a single
ECN-Echo flag, and the sender reacts at most once per window of data, about once per
RTT~\cite{rfc3168}. A
QUIC ACK frame can carry cumulative counts of the packets received with each ECN
codepoint, ECT(0), ECT(1) and CE, separately for each packet number
space~\cite{rfc9000}. Packets marked CE should be acknowledged immediately, and each
endpoint validates the counts on every path and disables ECN on a path that
corrupts them. The counts tell a sender which fraction of its packets
was marked in each round. This supports a response proportional to the extent of
congestion and an exit from startup before any packet is lost. TCP can now also
receive more than one ECN signal per RTT, with Accurate ECN~\cite{rfc9768}.

The main limitation of ECN in QUIC is deployment,
and the main obstacle are the servers. In 2023, only 0.22\% of QUIC
domains passed ECN validation~\cite{sander2023ecn}. Only 5.6\% of the domains
reported the ECN codepoints of received packets back to the client at all. Large
providers such as Cloudflare and Fastly did not, although no interference was visible
on the path, and Cloudflare supports ECN over TCP, which points to QUIC stacks that
do not implement ECN. Of the domains that did report the codepoints, 96\% then failed
validation. In 65\% of these failures, the server reported fewer ECN-marked packets
than were sent, which the authors attribute mostly to bugs in the QUIC stacks. In
31\%, it reported ECT(1) instead of the ECT(0) codepoint that was sent, caused either
by routers that rewrite the codepoint or by stacks that confuse the
two. QUIC is strict about the ECN counts: it compares the counts with what it sent and disables ECN on any mismatch.

\paragraph{Control over acknowledgment frequency.} In the base protocol, as in TCP,
the receiver decides when to acknowledge. It should send an ACK after receiving at
least two ack-eliciting packets, and it must not delay an acknowledgment by more than
its \texttt{max\_ack\_delay}. It must acknowledge Initial and Handshake packets
immediately, and it should acknowledge immediately when packets arrive out of order or
are marked CE~\cite{rfc9000}. Control by the sender comes from the QUIC
acknowledgment frequency extension~\cite{iyengar2026ackfrequency}. Its
ACK\_FREQUENCY frame lets the sender set how many ack-eliciting packets the receiver
may receive before acknowledging, the maximum acknowledgment delay it requests, and
the reordering threshold. Its IMMEDIATE\_ACK frame asks the receiver to acknowledge
at once.

Encryption also protects the acknowledgments themselves. TCP acknowledgments are
visible to the network, and on wireless and satellite links middleboxes thin or
filter them~\cite{custura2023reducing}. QUIC carries acknowledgments inside its
encrypted payload, so ACK manipulation is impossible~\cite{rfc9312}. 

\section{Flow sizes in Internet traffic}\label{sec:bg-traffic}

Measurements taken on different networks over more than two decades agree on two
points: most flows carry few bytes, and most bytes are carried by a small number of
large flows. Table~\ref{tab:bg-flow-sizes} summarizes the studies used in this
section.

\begin{table}[H]
  \centering
  \small
  \caption{Measurements of flow sizes.}
  \label{tab:bg-flow-sizes}
  \begin{tabular}{>{\raggedright\arraybackslash}p{0.19\linewidth}
                  >{\raggedright\arraybackslash}p{0.29\linewidth}
                  >{\raggedright\arraybackslash}p{0.40\linewidth}}
    \toprule
    Study & Vantage point and period & Key result \\
    \midrule
    Zhang et al.~\cite{zhang2002characteristics} & 8 packet traces of access and
      peering links, 2002 & Over 84\% of flows under 100\,KB; large, fast flows
      carry at least 60\% of bytes \\
    Bauer et al.~\cite{bauer2021evolution} & CAIDA Tier-1 backbone links,
      2008--2019, over 2.5 billion TCP flows & About 85\% of flows carry
      100\,B--10\,kB; largest 1\% carry about 89\% of bytes \\
    Jurkiewicz et al.~\cite{jurkiewicz2021flow} & University uplink, 30 days in
      2015, 4 billion flows & 92\% of TCP flows carry at most 16\,KiB and 1.2\% of
      TCP bytes \\
    Luxemburk et al.~\cite{luxemburk2023cesnet} & National research network
      backbone, 2022, 153 million QUIC flows & 47.5\% of QUIC flows carry under
      10\,KB; 98.0\% under 1\,MB \\
    Cloudflare~\cite{cloudflare2025connections} & CDN servers, 1\% sample of TCP
      connections, 2025 & Median connection receives 4.8--6\,KB from the server;
      90th percentile 51--146\,KB \\
    \bottomrule
  \end{tabular}
\end{table}

\paragraph{Flow sizes.} Zhang et al.\ found that more than 84\% of the flows in each
of their eight packet traces were smaller than 100\,KB, while the few flows that
were both large and fast carried at least 60\% of the
bytes~\cite{zhang2002characteristics}. More recent data show that the pattern
has persisted. In more than 2.5 billion TCP flows from CAIDA traces of Tier-1
backbone links taken between 2008 and 2019, Bauer et al.\ found that about 85\% of
the flows carry between 100\,B and 10\,kB of payload, while the largest 1\% of flows
carry about 89\% of the bytes~\cite{bauer2021evolution}. In 4 billion flows
collected over 30 days at a university uplink, Jurkiewicz et al.\ found that 92\% of
TCP flows carry at most 16\,KiB and together account for only 1.2\% of the TCP
bytes, while flows larger than 1\,MiB are 0.66\% of the TCP flows but carry 91\% of
the TCP bytes~\cite{jurkiewicz2021flow}.

\paragraph{QUIC traffic.} The flow sizes of QUIC traffic have been studied far less
than those of TCP, although QUIC already carries a large share of web traffic: in
2025, HTTP/3, which runs over QUIC, accounted for 21\% of the requests to
Cloudflare~\cite{cloudflare2025radar}. The closest source that studies QUIC traffic is the CESNET-QUIC22
dataset~\cite{luxemburk2023cesnet}, which contains 153 million QUIC flows observed on
the backbone of CESNET2, the Czech national research and education network, between
31 October and 27 November 2022. Of these flows, 47.5\% carry less than 10\,KB in
both directions together, 98.0\% carry less than 1\,MB, and 59.2\% consist of fewer
than 30 packets. QUIC flows therefore follow similar patterns as TCP flows: most of
them carry few bytes.

\paragraph{Web traffic} The web shows why these flow sizes matter for
congestion control. A median web page is large, about 2.2--2.4\,MB, but it is fetched
as 72--77 requests, mostly for small objects: the median image is 8\,KB and the
median JavaScript file is smaller than 5\,KB~\cite{httparchive2025pageweight}. The
connections that carry these requests are small as well. In a 1\% sample of the TCP
connections to Cloudflare's servers in October 2025, which excludes HTTP/3, the median
connection received only 4.8\,KB from the server over HTTP/1.x and 6\,KB over HTTP/2,
the 90th percentiles were 51\,KB and 146\,KB, and the median connection carried a
single request~\cite{cloudflare2025connections}.

\paragraph{Congestion control for web traffic} Most web transfers end while the
sender is still in slow start. Cardwell et al.\ observed that most TCP flows are so
short that they often spend their entire lifetime in slow start, without suffering a
single loss~\cite{cardwell2000modeling}. Dukkipati et al.\ report the same for web
connections: the majority finish before exiting slow start, so their latency is
dominated by the number of RTTs that slow start takes~\cite{dukkipati2010argument}.

Compared with the windows that congestion control works with, these transfers are
small. At least half of the connections to Cloudflare's servers fit into the initial
window, and 90\% of the HTTP/1.x connections are smaller than the median BDP of
Cloudflare's connections, about 77\,KB~\cite{cloudflare2025connections}. A flow that
fits into the initial window completes one RTT after the handshake, whatever the CCA
does, and a flow smaller than the BDP ends before its window can fill the path.

For a flow larger than the initial window, the number of rounds depends on how fast
the window grows. A sender that starts with an initial window of $\mathit{IW}$ bytes
and doubles its window every RTT needs
\begin{equation}
  n = \left\lceil \log_2\!\left(\frac{S}{\mathit{IW}} + 1\right) \right\rceil
  \label{eq:bg-ss-rounds}
\end{equation}
rounds to send $S$ bytes without loss. A 100\,KB transfer with a 15\,KB initial
window therefore needs three rounds, even on a path whose BDP would let it finish in
one.

Web transfers therefore fall into three groups. The smallest requests fit into the
initial window, and no CCA can make them faster: they complete one RTT after the
handshake, so only fewer handshake round trips or a shorter RTT reduce their latency.
Large transfers, much larger than the BDP, spend only a small part of their lifetime
in startup, so their FCT depends mainly on the throughput in steady state. Between
them are medium-sized flows, which spend most or all of their lifetime in startup. For these flows, reaching
the BDP in as few rounds as possible directly reduces the number of RTTs, and
therefore the FCT, provided that the sender does not overshoot the BDP and cause
losses that take further round trips to recover.


\paragraph{Consequences of delay} On the web, these RTTs
add directly to the delay a user sees, and that delay has a measurable cost. In a
controlled experiment at Bing, which slowed the search results page by 100 or
250\,ms for a subset of users, every 100\,ms of speedup increased revenue by
0.6\%; Amazon similarly reported that a slowdown of 100\,ms decreased its sales by
1\%~\cite{kohavi2014seven}. At Google, adding 100 to 400\,ms of delay to the search
results reduced the number of daily searches per user by 0.2 to 0.6\%, and after a
delay of 400\,ms, the effect persisted even once the delay was
removed~\cite{brutlag2009speed}.
